% ---------------------------------------------------------------------------------------------------------
% 縮約モデルと軌道最適化
% ---------------------------------------------------------------------------------------------------------
\section{縮約モデルと軌道最適化}
\label{sec:model}
\label{sec:opt}
\subsection{平面 5 リンク人体}
人体を，前腕＋手，上腕，体幹＋頭部，大腿，下腿の 5 本の剛体リンクで表す（図~\ref{fig:schematic} 左）．左右の腕と脚はそれぞれ 1 本にまとめた．
各リンクの質量比・重心位置・回転半径は de~Leva~\cite{deleva1996} の値を用い，体重 $m$ と身長 $H_b$（基準 1.75\,m，
全リンク長を $H_b/1.75$ 倍する）に合わせて構成した．
両手の掛かり線の位置 $\bm p_h$ とリンク角 $\theta_1,\dots,\theta_5$ を並べた一般化座標 $\bm q=(\bm p_h,\theta_1,\dots,\theta_5)$ について，運動方程式は
\begin{equation}
\bm M(\bm q)\ddot{\bm q}+\bm h(\bm q,\dot{\bm q})=\bm B\bm\tau+\bm J_c^\top\bm R
\label{eq:eom}
\end{equation}
である．$\bm M$ は質量行列，$\bm h$ は重力と速度に依る項，$\bm\tau=(\tau_1,\dots,\tau_4)$ は肘・肩・股・膝の関節トルク，
$\bm B$ はトルクを一般化力に写す行列，$\bm J_c$ は手の位置のヤコビアン，$\bm R$ は突起反力である．
式は CasADi~\cite{andersson2019} で記号的に導出した．
把持中は $\bm p_h$ が突起先端に一致して動くので，$\bm R$ は式~\eqref{eq:eom} の手の 2 行から拘束力として決まる．
モデルの定数を表~\ref{tab:constants} に示す．
関節トルクには上限 $|\tau_j|\le\taucap_j$，関節角には可動域，関節角速度には上限 10\,rad/s，
指令トルクの変化率には上限 $10\,\taucap_j$/s を課した．
\begin{table*}[t]
\centering\scriptsize
\caption{モデルの定数．トルク上限は左右の合計．可動域は隣り合うリンクの角度差で，A に向いているときの値である（B に向くと符号が反転する）．}
\label{tab:constants}
\resizebox{\textwidth}{!}{%
\begin{tabular}{llll}
\toprule
項目 & 値 & 項目 & 値 \\
\midrule
関節トルク上限 $\taucap_1,\dots,\taucap_4$（肘，肩，股，膝） & 120，160，300，240\,N\,m & 引っ掛かり係数 $\mu_{\mathrm{out}}$，摩擦係数 $\mu_{\mathrm{in}}$ & 1.0，2.0 \\
肘の可動域 $\theta_1-\theta_2$ & 0〜149 度（0 が伸展） & 突起の奥行き・高さ・本体前面 & 3\,cm，5\,cm，15\,cm \\
肩の可動域 $\theta_3-\theta_2$ & $-241$〜34 度（過伸展 34 度まで） & 先端上縁の丸み & 半径 5\,mm \\
股の可動域 $\theta_4-\theta_3$ & $-120$〜20 度 & 壁との隙間（肘，肩，股，膝，足先） & 4，12，12，7，10\,cm \\
膝の可動域 $\theta_5-\theta_4$ & 0〜149 度（0 が伸展） & 前腕上の干渉点（掛かり線から 8，16\,cm）の隙間 & 1，2\,cm \\
関節角速度の上限 & 10\,rad/s & 捕捉時の相対速度の上限 $\|\bm v_{\mathrm{rel}}\|$ & 4\,m/s \\
指令トルクの変化率の上限 & $10\,\taucap_j$/s & 捕捉時に壁から離れる向きの速度の上限 & 0.5\,m/s \\
把持力の表示単位 $\Fref$ & 1300\,N & 身長の基準 $H_b$ & 1.75\,m \\
\bottomrule
\end{tabular}}
\end{table*}

\textbf{前腕干渉．}把持中の前腕は突起の下の壁面に沿って垂れ下がる．そこで，掛かり線から 8\,cm と 16\,cm の前腕上の点が
壁面より手前（余裕 1\,cm，2\,cm）にあるという制約を全相で課した．
初期の実装にはこの制約がなく，最適化器は前腕を壁板に食い込ませて振り上げる非物理的な解を返していた
（3 次元モデルの構築中に発見した）．本稿の結果はすべて，この制約を入れたモデルで計算し直したものである．

\subsection{相の構成と捕捉モデル}
運動を，待機（A を持って静止），振り（A に拘束），飛行（どこにも触れない），捕捉，保持（B に拘束して $T_{\mathrm{hold}}=2.0$\,s）の
相に分ける多相問題として定式化した．各相の所要時間 $d_w$，$d_s$，$d_f$ は最適化の変数である．
飛行中は関節トルクだけが働き，体の向きの反転（ひねり）は矢状面のモデルでは自由に行えるものとした
（関節可動域は，振りでは A 向き，保持では B 向き，飛行ではその両方を許す）．

捕捉は 2 通りにモデル化した．
\textbf{剛体衝突モデル}は，掛かり線が B の先端に一致し，手が把持面へ向かって動いている
（鉛直の相対速度が下向き，壁から離れる向きの速度が 0.5\,m/s 以下，相対速度の大きさが 4\,m/s 以下）ときに，
手が瞬時に突起に固定されるとする．このとき働く衝撃力積 $\bm\Lambda$ は
$\bm M(\dot{\bm q}^+-\dot{\bm q}^-)=\bm J_c^\top\bm\Lambda$（$\dot{\bm q}^\mp$ は衝突の前後の速度）から決まる．
捕捉の負荷は，力積を時間 $\Delta_c=50$\,ms に均した力と，捕捉直後の反力 $\bm R_H(0)$ を足して
\begin{equation}
U_{\mathrm{catch}}=\frac{\|\bm\Lambda\|/\Delta_c+\|\bm R_H(0)\|}{\Fref}
\end{equation}
と見積もる．
\textbf{順応捕捉モデル}は，指を剛性 $K_c=40$\,kN/m，減衰 $D_c=1.5$\,kN\,s/m のばね・ダンパとして扱い，
捕捉の 0.1\,s を独立した相として最適化の内部に置く．両者の差は第~\ref{sec:interventions}節で比べる．

突起反力には，片側支持 $R_y\ge0$ と摩擦・引っ掛かり錐
\begin{equation}
-\mu_{\mathrm{out}}R_y\le sR_x\le\mu_{\mathrm{in}}R_y
\end{equation}
を課した（図~\ref{fig:schematic} 右；$s=+1$ は A，$s=-1$ は B で，壁から離れる向きを負にとる）．
壁および突起との干渉は，体の各点が「壁の前面より手前にある」か「壁の下端より下にある」のいずれかを満たすという
条件を，滑らかな最大値関数で表して課した．

\subsection{目的関数}
主目的は必要把持容量の最小化である．努力指標を
\begin{equation}
\Eeff=\frac{1}{t_{\mathrm{ref}}}\int\Bigl\{\frac14\sum_{j=1}^{4}u_j^2+U^2\Bigr\}\,dt
\label{eq:effort}
\end{equation}
とし（$u_j=\tau_j/\taucap_j$，$t_{\mathrm{ref}}=1$\,s は無次元化のための時間），
\begin{equation}
J=w_U\,\Upk+w_E\,\Eeff,\qquad w_U=10,\ w_E=1
\label{eq:objective}
\end{equation}
を最小化した．$\Upk$ は全時刻で $U(t)\le\Upk$ を満たす補助変数として扱う．
重みの選び方が結論を変えないことは第~\ref{sec:objective}節で確かめる．

\subsection{数値解法}
各相を圧縮型 Hermite--Simpson 法~\cite{hargraves1987,kelly2017} で離散化し（振り $N_S=150$，飛行 $N_F=30$，保持 $N_H=60$ 区間），
IPOPT~\cite{wachter2006} で解いた．
問題は非凸で，初期解によって別の局所解に落ちる．
そこで各条件を，振幅・周期・立ち上がりの異なる 6 種類の振りの初期解から解いた．
さらに，装置周期 $T$，体重 $m$，離手位相 $\phi_\ell$ の格子上で隣り合う条件の解を初期値にして解き直す\textbf{連続法}を，
改善がなくなるまで繰り返した．
連続法が必要な理由は第~\ref{sec:grid}節で示す．
飛行中は重心が放物線を描き角運動量が保存されるという条件を許容幅つきの制約として加え，
得られた解はすべて刻み 1\,ms の 4 次 Runge--Kutta 法で積分し直して一致を確かめた．

\subsection{双対変数と感度}
最適化の副産物として，各制約のラグランジュ乗数が得られる．
関節トルク上限や錐の係数などを最適化問題のパラメータ $p$ として定式化しておくと，
包絡線定理~\cite{fiacco1976,buskens2001} により，最適値のパラメータについての勾配 $\partial J^*/\partial\ln p$ が乗数から計算できる．
これは「その能力が 1\,\%増えたとき最適値がどれだけ下がるか」を表す影の価格である．
また，$U(t)\le\Upk$ の乗数を相ごとに合計すると，$\Upk$ を決めているのがどの相かが分かる．
ミニマックス解では複数の時刻が同時に $\Upk$ に達するので，到達しているかどうかだけでは支配的な相を判定できない．
乗数は，その相の制約を緩めたときに最適値が下がる量を与えるので，この判定に使える．

\section{3 次元人体モデル}
\label{sec:spatial}
平面モデルは左右の手を 1 本にまとめるので，捕捉の衝撃が先に触れた手に集中する効果を表せない．
この点を調べるため，体幹 6 自由度と，左右独立の肩 2・肘 1・股 1・膝 1 自由度をもつ 14 自由度のひねりモデルを CasADi で構築し，
左右の手の離手・捕捉に時間差を許す多相最適化の枠組みを実装した．
また MuJoCo~\cite{todorov2012} 上に 31 関節の人体と有限容量の把持モデルを構築した．
図~\ref{fig:device3d} と以降の 3 次元の図は，この 31 関節モデルを平面モデルの姿勢に合わせて描いたものである．
3 次元モデルによる片手ごとの負荷の掃引は本稿の範囲外とし，以降の結果はすべて平面モデルによる．
